
\begin{equation}\tag{4.4.3}
 (f^*\underline{L}_1,\ldots,f^*\underline{L}_d;O_Y)
 =
 n(\underline{L}_1,\ldots,\underline{L}_d;\underline{O}_X).
\end{equation}
Rappelons la définition du degré \(n\) de \(f\) [3] : il est nul sauf dans le cas où \(\dim(X)=\dim(Y)=\dim(f(Y))\). Dans ce dernier cas on pose
\begin{equation}\tag{4.4.4}
 n=(\len_{\underline{O}_{X,x}}(\underline{O}_{Y,y}))/(\len_{\underline{O}_{X,x}}(\underline{O}_{X,x})),
\end{equation}
où \(x\) est le point générique de \(X\) et \(y\) le point générique de \(Y\).
\end{enumerate}

Dans la démonstration de (4.4.3) le cas \(\dim(X)=\dim(Y)\) est le seul non-trivial. Si \(\dim(f(Y))<d\), alors on a
\begin{equation}\tag{4.4.5}
 R^if_*(\underline{O}_Y)=0
\end{equation}
pour \(i>0\) dans un voisinage de \(x\). Mais la condition (4.4.5) est satisfaite aussi dans le cas \(\dim(f(Y))=d\), car alors la dimension de la fibre générique est nulle (EGA III 4.2.2 et EGA IV 5.6.6). Alors 1.1.2 donne les congruences
\[
 f_*(\cl_\bullet(\underline{O}_Y))
 \equiv \cl_\bullet(f_*(\underline{O}_Y)) \equiv
 \len_{\underline{O}_{X,x}}(\underline{O}_{Y,y})
 \cl_\bullet(\underline{O}_{X_{\mathrm{red}}})
\]
modulo \(\Filt_{d-1}(X)\). En vertu de (4.4.2) il en découle que
\[
 (f^*\underline{L}_1,\ldots,f^*\underline{L}_d;\underline{O}_Y)
 =
 \len_{\underline{O}_{X,x}}(\underline{O}_{Y,y})
 (\underline{L}_1,\ldots,\underline{L}_d;\underline{O}_{X_{\mathrm{red}}}).
\]
D'autre part \(2^\circ\) donne
\[
 (\underline{L}_1,\ldots,\underline{L}_d;\underline{O}_X)
 =
 \len_{\underline{O}_{X,x}}(\underline{O}_{X,x})
 (\underline{L}_1,\ldots,\underline{L}_d;\underline{O}_{X_{\mathrm{red}}}),
\]
d'où l'assertion.

